\SourceLeaf{438}%
\[
\begin{aligned}
\varrho d\xi&=Xds\\
\varrho d\eta&=Yds\\
\varrho d\zeta&=Zds
\end{aligned}
\]
und wir haben hier statt $4$ gekrümmter Linien auf der Oberfläche der Hilfskugel nur $3$ zu betrachten. Jedes Element der zweiten Linie ist also als im Großkreis $L\lambda$ liegend anzusehen, und der positive oder negative Wert von $\varrho$ bezieht sich auf die Konkavität oder Konvexität der Kurve in Richtung der Normalen.

\Needspace{5\baselineskip}
\begin{center}14.\end{center}
\noindent\hspace*{1.2em}Wir wollen nun den sphärischen Winkel auf der Hilfskugel untersuchen, den der Großkreis von $L$ nach $\lambda$ mit demjenigen bildet, der von $L$ zu einem der festen Punkte $(1),(2),(3)$, z.\,B. nach $(3)$, verläuft. Um dies genau festzulegen, messen wir von $L(3)$ nach $L\lambda$ in demjenigen Umlaufsinn, in dem $(1),(2)$ und $(3)$ liegen. Nennen wir $\varphi$ diesen Winkel, so gilt nach dem Satz in Art. $7$
\[
\sin L(3).\sin L\lambda.\sin\varphi=Y\xi-X\eta
\]
oder, da $L\lambda=90^{0}$ und
\[
\sin L(3)=\BareRoot{(XX+YY)}=\BareRoot{(1-ZZ)}:
\]
\[
\sin\varphi=\frac{Y\xi-X\eta}{\BareRoot{(XX+YY)}}.
\]
Ferner gilt
\[
\sin L(3).\sin L\lambda.\cos\varphi=\zeta
\]
oder
\[
\cos\varphi=\frac{\zeta}{\BareRoot{(XX+YY)}}
\]
und
\[
\operatorname{tang}\varphi=\frac{Y\xi-X\eta}{\zeta}=\frac{\zeta'}{\zeta}.
\]
Man erhält also
\[
d\varphi=\frac{\zeta Yd\xi-\zeta Xd\eta-(Y\xi-X\eta)d\zeta+\xi\zeta dY-\eta\zeta dX}{(Y\xi-X\eta)^{2}+\zeta\zeta}.
\]
